\begingroup
\renewcommand{\arraystretch}{1.45}
\fancyhead[C]{}
\providecommand{\TStart}{\hspace*{.5em}}\providecommand{\TCont}{\hspace*{1em}}\providecommand{\TD}{\hspace{.3em}.\hspace{.3em}.}
\clearpage\fancyhead[L]{\small INHALT.}\pdfbookmark[0]{Inhalt}{original-inhalt}\SourceLeaf{506}


\begin{center}
\textbf{INHALT.}\par
\medskip
\textbf{GAUSS WERKE BAND III.\quad ANALYSIS.}\par
\medskip
\rule{2em}{0.4pt}
\end{center}
\begingroup\small
\setlength{\tabcolsep}{0pt}
\setlength{\LTleft}{0pt}\setlength{\LTright}{0pt}
\begin{longtable}{@{}l@{\hspace{.01\linewidth}}p{.19\linewidth}@{\hspace{.01\linewidth}}p{.075\linewidth}@{\hspace{.005\linewidth}}p{.05\linewidth}@{}}
\multicolumn{4}{c}{\textbf{ALGEBRAISCHE FUNKTIONEN.}}\\*[.3em]
\multicolumn{4}{@{}l@{}}{\textit{Abhandlungen.}}\\*
\parbox[b]{.64\linewidth}{\raggedright Neuer Beweis des Satzes, dass jede rationale ganze algebraische Funktion einer Veränderlichen in reelle Faktoren ersten oder zweiten Grades zerlegt werden kann ~\hfill\mbox{.\kern.08em .\kern.08em .\kern.08em .\kern.08em .\kern.08em .\kern.08em .\kern.08em .\kern.08em .\kern.08em .\kern.08em .\kern.08em .\kern.08em .\kern.08em .\kern.08em .\kern.08em .\kern.08em .\kern.08em .\kern.08em .\kern.08em .\kern.08em .}\par} & 1799 & Seite & 1\\
\parbox[b]{.64\linewidth}{\raggedright Zweiter neuer Beweis des Satzes, dass jede rationale ganze algebraische Funktion einer Veränderlichen in reelle Faktoren ersten oder zweiten Grades zerlegt werden kann ~\hfill\mbox{.\kern.08em .\kern.08em .\kern.08em .\kern.08em .\kern.08em .\kern.08em .\kern.08em .\kern.08em .\kern.08em .\kern.08em .\kern.08em .\kern.08em .\kern.08em .\kern.08em .\kern.08em .\kern.08em .\kern.08em .\kern.08em .\kern.08em .}\par} & 1815 Dec.\TD & --- & 31\\
\parbox[b]{.64\linewidth}{\raggedright Dritter Beweis des Satzes über die Zerlegbarkeit ganzer algebraischer Funktionen in reelle Faktoren. Ergänzung zur vorangehenden Abhandlung\par} & 1816 Jan.\TD & --- & 57\\
\parbox[b]{.64\linewidth}{\raggedright Beweis eines algebraischen Lehrsatzes ~\hfill\mbox{.\kern.08em .\kern.08em .\kern.08em .\kern.08em .\kern.08em .\kern.08em .\kern.08em .\kern.08em .\kern.08em .\kern.08em .\kern.08em .\kern.08em .}\par} & 1828 & --- & 65\\
\parbox[b]{.64\linewidth}{\raggedright Beiträge zur Theorie der algebraischen Gleichungen ~\hfill\mbox{.\kern.08em .\kern.08em .\kern.08em .\kern.08em .\kern.08em .\kern.08em .\kern.08em .}\par} & 1849 Juli\TD & --- & 71\\
\multicolumn{4}{@{}l@{}}{\textit{Anzeigen eigener Abhandlungen.}}\\*
\parbox[b]{.64\linewidth}{\raggedright Zweiter neuer Beweis des Satzes, dass jede rationale ganze algebraische Funktion einer Veränderlichen in reelle Faktoren ersten oder zweiten Grades zerlegt werden kann ~\hfill\mbox{.\kern.08em .\kern.08em .\kern.08em .\kern.08em .\kern.08em .\kern.08em .\kern.08em .\kern.08em .\kern.08em .\kern.08em .\kern.08em .\kern.08em .\kern.08em .\kern.08em .\kern.08em .\kern.08em .\kern.08em .\kern.08em .}\par} & 1815 Dec.\TD & --- & 105\\
\parbox[b]{.64\linewidth}{\raggedright Dritter Beweis des Satzes über die Zerlegbarkeit ganzer algebraischer Funktionen in reelle Faktoren ~\hfill\mbox{.\kern.08em .\kern.08em .\kern.08em .\kern.08em .\kern.08em .\kern.08em .\kern.08em .\kern.08em .\kern.08em .\kern.08em .\kern.08em .\kern.08em .\kern.08em .}\par} & 1816 März\TD & --- & 107\\
\parbox[b]{.64\linewidth}{\raggedright Beiträge zur Theorie der algebraischen Gleichungen ~\hfill\mbox{.\kern.08em .\kern.08em .\kern.08em .\kern.08em .\kern.08em .\kern.08em .\kern.08em .}\par} & 1849 Juli\TD & --- & 113\\
\multicolumn{4}{@{}l@{}}{\textit{Anzeigen fremder Schriften.}}\\*
\parbox[b]{.64\linewidth}{\raggedright \textsc{Coytier.} Mathematische Untersuchungen ~\hfill\mbox{.\kern.08em .\kern.08em .\kern.08em .\kern.08em .\kern.08em .\kern.08em .\kern.08em .\kern.08em .\kern.08em .\kern.08em .\kern.08em .\kern.08em .}\par} & 1813 März\TD & --- & 116\\
\parbox[b]{.64\linewidth}{\raggedright \textsc{Fourier.} Analysis bestimmter Gleichungen ~\hfill\mbox{.\kern.08em .\kern.08em .\kern.08em .\kern.08em .\kern.08em .\kern.08em .\kern.08em .\kern.08em .\kern.08em .\kern.08em .}\par} & 1833 Febr.\TD & --- & 119\\[1em]
\multicolumn{4}{c}{\textbf{GAUSS’ REIHE.}}\\*
\multicolumn{4}{@{}l@{}}{\textit{Abhandlungen.}}\\*
\multicolumn{4}{@{}l@{}}{\TStart Allgemeine Untersuchungen über die unendliche Reihe}\\[.3em]
\noalign{\nobreak}
\multicolumn{4}{@{}l@{}}{\TCont $\displaystyle\begin{gathered}
1+\frac{\alpha.\beta}{1.\gamma}x
+\frac{\alpha(\alpha+1)\beta(\beta+1)}{1.2.\gamma(\gamma+1)}xx
\\[9pt]+\frac{\alpha(\alpha+1)(\alpha+2)\beta(\beta+1)(\beta+2)}{1.2.3.\gamma(\gamma+1)(\gamma+2)}x^{3}
+\mathord{.}\,\mathord{.}\end{gathered}$}\\[.3em]
\noalign{\nobreak}
\parbox[b]{.64\linewidth}{\raggedright Erster Teil ~\hfill\mbox{.\kern.08em .\kern.08em .\kern.08em .\kern.08em .\kern.08em .\kern.08em .\kern.08em .\kern.08em .\kern.08em .\kern.08em .\kern.08em .\kern.08em .\kern.08em .\kern.08em .\kern.08em .\kern.08em .\kern.08em .\kern.08em .\kern.08em .\kern.08em .\kern.08em .}\par} & 1812 Jan.\TD & --- & 123\\
\end{longtable}
\endgroup


\SourceLeaf{507}


\begingroup\small
\setlength{\tabcolsep}{0pt}
\setlength{\LTleft}{0pt}\setlength{\LTright}{0pt}
\begin{longtable}{@{}l@{\hspace{.01\linewidth}}p{.19\linewidth}@{\hspace{.01\linewidth}}p{.075\linewidth}@{\hspace{.005\linewidth}}p{.05\linewidth}@{}}
\parbox[b]{.64\linewidth}{\raggedright Neue Methode, die Werte von Integralen durch Näherung zu finden ~\hfill\mbox{.\kern.08em .\kern.08em .}\par} & 1814 Sept.\TD & Seite & 163\\
\multicolumn{4}{@{}l@{}}{\textit{Anzeigen eigener Abhandlungen.}}\\*
\parbox[b]{.64\linewidth}{\raggedright Allgemeine Untersuchungen über die unendliche Reihe usw. Erster Teil. ~\hfill\mbox{.\kern.08em .\kern.08em .\kern.08em .\kern.08em .}\par} & 1812 Febr.\TD & --- & 197\\
\parbox[b]{.64\linewidth}{\raggedright Neue Methode, die Werte von Integralen durch Näherung zu finden ~\hfill\mbox{.\kern.08em .\kern.08em .}\par} & 1814 Sept.\TD & --- & 202\\
\multicolumn{4}{@{}l@{}}{\textit{Nachlass.}}\\*
\parbox[b]{.64\linewidth}{\raggedright Bestimmung unserer Reihe durch eine Differentialgleichung zweiter Ordnung ~\hfill\mbox{.}\par} & & --- & 207\\
\parbox[b]{.64\linewidth}{\raggedright Bemerkungen zu dieser Abhandlung ~\hfill\mbox{.\kern.08em .\kern.08em .\kern.08em .\kern.08em .\kern.08em .\kern.08em .\kern.08em .\kern.08em .\kern.08em .\kern.08em .\kern.08em .\kern.08em .\kern.08em .}\par} & & --- & 230\\[1em]
\multicolumn{4}{@{}p{.98\linewidth}@{}}{\centering\textbf{MITTEILUNGEN ÜBER VERSCHIEDENE SCHRIFTEN.}}\\*[.4em]
\parbox[b]{.64\linewidth}{\raggedright \textsc{Pfaff.} Allgemeine Methode zur vollständigen Integration partieller Differentialgleichungen und gewöhnlicher Differentialgleichungen, beide erster Ordnung, zwischen beliebig vielen Veränderlichen ~\hfill\mbox{.\kern.08em .\kern.08em .\kern.08em .\kern.08em .\kern.08em .\kern.08em .\kern.08em .\kern.08em .\kern.08em .\kern.08em .\kern.08em .}\par} & 1815 Juli\TD & --- & 231\\
\parbox[b]{.64\linewidth}{\raggedright \textsc{Callet.} Logarithmische Tafeln ~\hfill\mbox{.\kern.08em .\kern.08em .\kern.08em .\kern.08em .\kern.08em .\kern.08em .\kern.08em .\kern.08em .\kern.08em .\kern.08em .\kern.08em .\kern.08em .\kern.08em .\kern.08em .\kern.08em .}\par} & 1802 Nov.\TD & --- & 241\\
\parbox[b]{.64\linewidth}{\raggedright \textsc{Prasse.} Logarithmische Tafeln für die Zahlen, Sinus und Tangenten ~\hfill\mbox{.\kern.08em .}\par} & 1811 Mai\TD & --- & 241\\
\parbox[b]{.64\linewidth}{\raggedright \textsc{Prasse.} Logarithmische Tafeln für Zahlen, Sinus und Tangenten\par} & 1814 Dec.\TD & --- & 243\\
\parbox[b]{.64\linewidth}{\raggedright \textsc{Gauss.} Tafeln zur bequemeren Berechnung des Logarithmen der Summe oder Differenz zweier Größen, welche selbst nur durch ihre Logarithmen gegeben sind ~\hfill\mbox{.\kern.08em .\kern.08em .\kern.08em .\kern.08em .\kern.08em .\kern.08em .\kern.08em .\kern.08em .\kern.08em .\kern.08em .\kern.08em .\kern.08em .\kern.08em .\kern.08em .\kern.08em .\kern.08em .\kern.08em .\kern.08em .\kern.08em .\kern.08em .\kern.08em .}\par} & 1812 Nov.\TD & --- & 244\\
\parbox[b]{.64\linewidth}{\raggedright \textsc{Pasquich.} Abgekürzte Logarithmisch-Trigonometrische Tafeln ~\hfill\mbox{.\kern.08em .\kern.08em .\kern.08em .\kern.08em .}\par} & 1817 Oct.\TD & --- & 246\\
\parbox[b]{.64\linewidth}{\raggedright \relax[\textsc{Matthiessen.}] Tafel zur bequemeren Berechnung des Logarithmen der Summe oder Differenz zweier Größen, welche selbst nur durch ihre Logarithmen gegeben sind ~\hfill\mbox{.\kern.08em .\kern.08em .\kern.08em .\kern.08em .\kern.08em .\kern.08em .\kern.08em .\kern.08em .\kern.08em .\kern.08em .\kern.08em .\kern.08em .\kern.08em .\kern.08em .\kern.08em .\kern.08em .\kern.08em .\kern.08em .}\par} & 1819 Jan.\TD & --- & 250\\
\parbox[b]{.64\linewidth}{\raggedright Logarithmentafeln mit sechs Ziffern ~\hfill\mbox{.\kern.08em .\kern.08em .\kern.08em .\kern.08em .\kern.08em .\kern.08em .\kern.08em .\kern.08em .\kern.08em .\kern.08em .\kern.08em .\kern.08em .\kern.08em .\kern.08em .}\par} & 1822 Dec.\TD & --- & 252\\
\parbox[b]{.64\linewidth}{\raggedright \textsc{Babbage.} Tafel der Logarithmen der natürlichen Zahlen ~\hfill\mbox{.\kern.08em .\kern.08em .\kern.08em .\kern.08em .\kern.08em .\kern.08em .\kern.08em .}\par} & 1828 Jan.\TD & --- & 253\\
\parbox[b]{.64\linewidth}{\raggedright \textsc{Hassler.} Logarithmische und trigonometrische Tafeln ~\hfill\mbox{.\kern.08em .\kern.08em .\kern.08em .\kern.08em .\kern.08em .\kern.08em .\kern.08em .\kern.08em .}\par} & 1831 März\TD & --- & 255\\
\parbox[b]{.64\linewidth}{\raggedright \textsc{Vega} und \textsc{Hülsse.} Sammlung mathematischer Tafeln ~\hfill\mbox{.\kern.08em .\kern.08em .\kern.08em .\kern.08em .\kern.08em .}\par} & 1840 & --- & 255\\
\parbox[b]{.64\linewidth}{\raggedright \textsc{Vega.} Schatz der Logarithmen ~\hfill\mbox{.\kern.08em .\kern.08em .\kern.08em .\kern.08em .\kern.08em .\kern.08em .\kern.08em .\kern.08em .\kern.08em .\kern.08em .\kern.08em .\kern.08em .\kern.08em .\kern.08em .}\par} & 1851 Mai\TD & --- & 257\\[1em]
\multicolumn{4}{c}{\textbf{INTERPOLATION.}}\\*
\multicolumn{4}{@{}l@{}}{\textit{Nachlass.}}\\*
\parbox[b]{.64\linewidth}{\raggedright Theorie der Interpolation nach einer neuen Methode behandelt ~\hfill\mbox{.\kern.08em .\kern.08em .\kern.08em .\kern.08em .\kern.08em .\kern.08em .\kern.08em .\kern.08em .\kern.08em .\kern.08em .\kern.08em .\kern.08em .\kern.08em .\kern.08em .\kern.08em .}\par} & & --- & 265\\
\parbox[b]{.64\linewidth}{\raggedright Bemerkungen zu dieser Abhandlung ~\hfill\mbox{.\kern.08em .\kern.08em .\kern.08em .\kern.08em .\kern.08em .\kern.08em .\kern.08em .\kern.08em .\kern.08em .\kern.08em .\kern.08em .\kern.08em .\kern.08em .\kern.08em .\kern.08em .\kern.08em .\kern.08em .\kern.08em .}\par} & & --- & 328\\[1em]
\multicolumn{4}{c}{\textbf{ELLIPTISCHE FUNKTIONEN.}}\\*
\multicolumn{4}{@{}l@{}}{\textit{Abhandlung.}}\\*
\parbox[b]{.64\linewidth}{\raggedright Bestimmung der Anziehung,. die ein Planet auf einen beliebigen Punkt gegebener Lage ausüben würde, wenn seine Masse über die ganze Bahn im Verhältnis der Zeit, in der die einzelnen Teile durchlaufen werden, gleichmäßig verteilt wäre ~\hfill\mbox{.\kern.08em .\kern.08em .\kern.08em .\kern.08em .\kern.08em .\kern.08em .\kern.08em .\kern.08em .}\par} & 1818 Jan.\TD & --- & 331\\
\parbox[b]{.64\linewidth}{\raggedright \textit{Anzeige der vorhergehenden Abhandlung} ~\hfill\mbox{.\kern.08em .\kern.08em .\kern.08em .\kern.08em .\kern.08em .\kern.08em .\kern.08em .\kern.08em .\kern.08em .\kern.08em .\kern.08em .\kern.08em .\kern.08em .\kern.08em .}\par} & 1818 Febr.\TD & --- & 357\\
\multicolumn{4}{@{}l@{}}{\textit{Nachlass.}}\\*
\multicolumn{4}{@{}p{.98\linewidth}@{}}{\TStart Über den Ursprung und die allgemeinen Eigenschaften arithmetisch-geometrischer Mittelzahlen.}\\
\parbox[b]{.64\linewidth}{\raggedright Art.\ 1\textemdash 8.~\hfill\mbox{.\kern.08em .\kern.08em .\kern.08em .\kern.08em .\kern.08em .\kern.08em .\kern.08em .\kern.08em .\kern.08em .\kern.08em .\kern.08em .\kern.08em .\kern.08em .\kern.08em .\kern.08em .\kern.08em .\kern.08em .\kern.08em .\kern.08em .\kern.08em .\kern.08em .\kern.08em .\kern.08em .\kern.08em .\kern.08em .\kern.08em .\kern.08em .\kern.08em .}\par} & & --- & 361\\
\end{longtable}
\endgroup

\SourceLeaf{508}


\begingroup\small
\setlength{\tabcolsep}{0pt}
\setlength{\LTleft}{0pt}\setlength{\LTright}{0pt}
\begin{longtable}{@{}l@{\hspace{.015\linewidth}}p{.075\linewidth}@{\hspace{.005\linewidth}}p{.05\linewidth}@{}}
\parbox[b]{.855\linewidth}{\raggedright Über die transzendenten Funktionen, die aus der Differentiation arithmetisch-geometrischer Mittel entstehen. Art.\ 9\textemdash 11.~\hfill\mbox{.\kern.08em .\kern.08em .\kern.08em .\kern.08em .\kern.08em .\kern.08em .\kern.08em .\kern.08em .\kern.08em .\kern.08em .\kern.08em .\kern.08em .\kern.08em .\kern.08em .\kern.08em .\kern.08em .\kern.08em .\kern.08em .\kern.08em .\kern.08em .\kern.08em .}\par} & Seite & 372\\
\parbox[b]{.855\linewidth}{\raggedright Fortsetzung der Untersuchungen über das arithmetisch-geometrische Mittel. Art.\ 12\textemdash 26.~\hfill\mbox{.}\par} & --- & 375\\
\parbox[b]{.855\linewidth}{\raggedright Tafel. Arithmetisch-geometrische Mittel zwischen der Einheit und den Sinuswerten der einzelnen halben Grade ~\hfill\mbox{.\kern.08em .}\par} & --- & 403\\
\parbox[b]{.855\linewidth}{\raggedright Lemniskatische Funktionen (I.) dargestellt durch Potenzreihen ~\hfill\mbox{.\kern.08em .\kern.08em .\kern.08em .\kern.08em .\kern.08em .\kern.08em .\kern.08em .\kern.08em .\kern.08em .\kern.08em .}\par} & --- & 404\\
\parbox[b]{.855\linewidth}{\raggedright Lemniskatische Funktionen (II.) dargestellt durch unendliche Produkte und durch trigonometrische Reihen ~\hfill\mbox{.\kern.08em .\kern.08em .\kern.08em .\kern.08em .\kern.08em .\kern.08em .\kern.08em .\kern.08em .\kern.08em .\kern.08em .\kern.08em .\kern.08em .\kern.08em .\kern.08em .\kern.08em .\kern.08em .\kern.08em .\kern.08em .\kern.08em .\kern.08em .\kern.08em .\kern.08em .\kern.08em .\kern.08em .\kern.08em .\kern.08em .}\par} & --- & 413\\
\multicolumn{3}{@{}l@{}}{Zur Theorie der neuen Transzendenten}\\*
\parbox[b]{.855\linewidth}{\raggedright I.\quad Anfänge der der Untersuchungen über die neuen Transzendenten ~\hfill\mbox{.\kern.08em .\kern.08em .\kern.08em .\kern.08em .\kern.08em .\kern.08em .\kern.08em .}\par} & --- & 433\\
\parbox[b]{.855\linewidth}{\raggedright II.\quad Zur Theorie der transzendenten Funktionen gehörig ~\hfill\mbox{.\kern.08em .\kern.08em .\kern.08em .\kern.08em .\kern.08em .\kern.08em .\kern.08em .\kern.08em .\kern.08em .\kern.08em .\kern.08em .\kern.08em .}\par} & --- & 436\\
\parbox[b]{.855\linewidth}{\raggedright III.\quad Neue Anordnung der Theoreme in Bezug auf diejenigen Reihen und unendlichen Produkte, welche zu der Theorie der arithmetisch-geometrischen Mittel gehören ~\hfill\mbox{.\kern.08em .\kern.08em .\kern.08em .}\par} & --- & 446\\
\parbox[b]{.855\linewidth}{\raggedright IV.\quad Hundert Theoreme über die neuen Transzendenten ~\hfill\mbox{.\kern.08em .\kern.08em .\kern.08em .\kern.08em .\kern.08em .\kern.08em .\kern.08em .\kern.08em .\kern.08em .\kern.08em .\kern.08em .\kern.08em .}\par} & --- & 461\\
\parbox[b]{.855\linewidth}{\raggedright V.\quad Allgemeine Theoreme ~\hfill\mbox{.\kern.08em .\kern.08em .\kern.08em .\kern.08em .\kern.08em .\kern.08em .\kern.08em .\kern.08em .\kern.08em .\kern.08em .\kern.08em .\kern.08em .\kern.08em .\kern.08em .\kern.08em .\kern.08em .\kern.08em .\kern.08em .\kern.08em .\kern.08em .\kern.08em .\kern.08em .\kern.08em .}\par} & --- & 470\\
\parbox[b]{.855\linewidth}{\raggedright Das wunderbare Pentagramm ~\hfill\mbox{.\kern.08em .\kern.08em .\kern.08em .\kern.08em .\kern.08em .\kern.08em .\kern.08em .\kern.08em .\kern.08em .\kern.08em .\kern.08em .\kern.08em .\kern.08em .\kern.08em .\kern.08em .\kern.08em .\kern.08em .\kern.08em .\kern.08em .\kern.08em .\kern.08em .\kern.08em .\kern.08em .\kern.08em .}\par} & --- & 481\\
\parbox[b]{.855\linewidth}{\raggedright \textit{Bemerkungen} ~\hfill\mbox{.\kern.08em .\kern.08em .\kern.08em .\kern.08em .\kern.08em .\kern.08em .\kern.08em .\kern.08em .\kern.08em .\kern.08em .\kern.08em .\kern.08em .\kern.08em .\kern.08em .\kern.08em .\kern.08em .\kern.08em .\kern.08em .\kern.08em .\kern.08em .\kern.08em .\kern.08em .\kern.08em .\kern.08em .\kern.08em .\kern.08em .\kern.08em .\kern.08em .\kern.08em .}\par} & --- & 491\\[.7em]
\multicolumn{3}{@{}l@{}}{\textit{Zwei Steindrucktafeln zu Seite 30 und 102.}}\\*
\end{longtable}
\endgroup
\begin{center}\rule{7em}{0.4pt}\end{center}

\clearpage\fancyhead[L]{}\pdfbookmark[0]{Druckvermerk}{original-colophon}\SourceLeaf{509}

\vspace*{0.24\textheight}
\begin{center}
{\large GÖTTINGEN,}\par
\medskip
\textbf{GEDRUCKT IN DER DIETERICHSCHEN\\UNIVERSITÄTS-DRUCKEREI}\par
\smallskip
\textsc{W. Fr. Kaestner.}
\end{center}
\clearpage
\endgroup
